% 由 main.tex 在 \appendix 之后以 % 由 main.tex 在 \appendix 之后以 % 由 main.tex 在 \appendix 之后以 % 由 main.tex 在 \appendix 之后以 \input{appendix-nlq-prompt} 引入
\section{自然语言提示词与调试说明补遗}
\label{app:prompt-supplement}

本附录给出提示词的完整分块内容与少样本示例，供读者对照正文实现或在答辩演示中复现。对应的提示词文件位于仓库 \texttt{backend/app/system\_prompt.txt}，与第~\ref{sec:detailed-implementation}~章第~\ref{subsec:nlq-impl}~节的说明一致；如二者有出入，以仓库当前文件为准。

\subsection{提示词分块}

\paragraph{（1）角色定义。}
告知模型它的任务是根据用户问句生成 Cypher 查询语句，且生成的语句必须是只读的，不允许包含任何写操作。

\paragraph{（2）Schema 描述。}
向模型说明当前图数据库的节点与关系结构：
\begin{itemize}
    \item 基础标签：\texttt{:item}，通用属性含 ID、Name、Type、Disc、GetWay，所有道具节点均带此标签
    \item 方块节点带 \texttt{:item:block} 双标签，在通用属性基础上增加 MineTool、ToolLevel
    \item 生物节点带 \texttt{:item:monster} 双标签，增加 Life、Attack
    \item 其他标签：\texttt{:group}（分组）、\texttt{:recipe}（配方）、\texttt{:smelt}（熔炼）、\texttt{:device}（设备）
    \item 关系类型：\texttt{IN\_GROUP}、\texttt{CONSUMES}、\texttt{PRODUCES}、\texttt{DROPS}、\texttt{TOOLMINEDROPS} 等 12 种
\end{itemize}

\paragraph{（3）多标签查询语法。}
\begin{itemize}
    \item 查询方块：\texttt{MATCH (n:item:block) WHERE n.Name CONTAINS '关键词'}
    \item 查询生物：\texttt{MATCH (n:item:monster) WHERE n.Name CONTAINS '关键词'}
    \item 查询纯道具：\texttt{MATCH (n:item) WHERE NOT n:block AND NOT n:monster}
\end{itemize}

\paragraph{（4）路径查询与返回形态。}
本块引导模型优先返回路径，即使用 \texttt{RETURN path} 而非仅返回孤立节点。返回路径时，后端可以直接从记录中提取节点与边，用于绘制关系图；若只返回孤立节点，边信息丢失，前端无法展示实体之间的连接关系。

关于查询方向：无向关系与有向关系在 Cypher 语义上并不等价。在探索性问句中使用无向关系可以扩大命中范围，但也可能匹配到额外的边。本文将此作为提示策略的一个选项，不作为普遍性结论。

\paragraph{（5）少样本示例（few-shot）补充。}
以下标题中“方块挖掘”“生物掉落”等来自本文案例数据，仅作模式演示。

\begin{lstlisting}[language=cypher, caption={方块挖掘（无方向路径）}]
-- 深积岩用什么工具挖
MATCH path = (n:item:block)-[:TOOLMINEDROPS]-(tool:item)
WHERE n.Name CONTAINS '深积岩'
RETURN path LIMIT 10
\end{lstlisting}

\begin{lstlisting}[language=cypher, caption={生物掉落（无方向路径）}]
-- 击败野人会掉落什么
MATCH path = (n:item:monster)-[:DROPS]-(drop:item)
WHERE n.Name CONTAINS '野人'
RETURN path LIMIT 20
\end{lstlisting}

\begin{lstlisting}[language=cypher, caption={配方原料（无方向路径）}]
-- 石剑的合成配方
MATCH path = (product:item)-[:PRODUCES]-(r:recipe)-[:CONSUMES]-(material:item)
WHERE product.Name CONTAINS '石剑'
RETURN path LIMIT 20
\end{lstlisting}

\paragraph{（6）输出格式。}
要求大语言模型返回 JSON，含 \texttt{cql} 与 \texttt{params} 字段。

\subsection{部署与调参中的经验性说明}

在表~\ref{tab:nlq-category} 的自然语言实例测试之外，实现过程中还有以下几点观察。优先使用 \texttt{RETURN path} 与无向关系，有利于表格与力导向关系图同时展示。多标签方案（\texttt{:item:block} 等）使查询路径比旧版桥接关系方案更短，更易于模型理解和生成。\texttt{Name CONTAINS} 适合含复合名称的道具查询。\texttt{\$param} 参数化与关键字黑名单是防止写操作误用的工程措施。以上均为实现过程中的经验性观察，不构成与公开基准数据集的对照评测。

\subsection{迁移场景下的提示词改写说明}

在将系统迁移到其他 Neo4j 知识图谱时，需要改写系统提示词。
现有提示词由两类内容组成，处理方式不同。

第一类是与领域无关的约束性说明，迁移时可以直接保留，无需修改。
具体包括：角色定义（只读 Cypher 生成器，不允许写操作）、输出格式要求（JSON，含 \texttt{cql} 与 \texttt{params} 字段）、路径优先返回策略（使用 \texttt{RETURN path} 代替孤立节点），以及无向关系的使用建议。

第二类是与迷你世界数据库绑定的描述，迁移时需要按目标数据库替换。
具体包括：Schema 描述（第（2）块，含节点标签列表、关系类型列表与各标签的属性键说明）和少样本示例（第（5）块，含示范问句及其对应的 Cypher 语句）。
将这两块的内容替换为目标数据库的对应描述后，其余部分保持不变，即可完成提示词的迁移。

少样本示例的数量与质量对迁移效果影响较大。
建议针对目标数据库各主要节点标签与关系类型各准备一至两个代表性问句，
覆盖单节点属性查询、路径查询和带过滤条件的查询三种形态，使模型对目标领域的典型 Cypher 模式有足够的参照。
 引入
\section{自然语言提示词与调试说明补遗}
\label{app:prompt-supplement}

本附录给出提示词的完整分块内容与少样本示例，供读者对照正文实现或在答辩演示中复现。对应的提示词文件位于仓库 \texttt{backend/app/system\_prompt.txt}，与第~\ref{sec:detailed-implementation}~章第~\ref{subsec:nlq-impl}~节的说明一致；如二者有出入，以仓库当前文件为准。

\subsection{提示词分块}

\paragraph{（1）角色定义。}
告知模型它的任务是根据用户问句生成 Cypher 查询语句，且生成的语句必须是只读的，不允许包含任何写操作。

\paragraph{（2）Schema 描述。}
向模型说明当前图数据库的节点与关系结构：
\begin{itemize}
    \item 基础标签：\texttt{:item}，通用属性含 ID、Name、Type、Disc、GetWay，所有道具节点均带此标签
    \item 方块节点带 \texttt{:item:block} 双标签，在通用属性基础上增加 MineTool、ToolLevel
    \item 生物节点带 \texttt{:item:monster} 双标签，增加 Life、Attack
    \item 其他标签：\texttt{:group}（分组）、\texttt{:recipe}（配方）、\texttt{:smelt}（熔炼）、\texttt{:device}（设备）
    \item 关系类型：\texttt{IN\_GROUP}、\texttt{CONSUMES}、\texttt{PRODUCES}、\texttt{DROPS}、\texttt{TOOLMINEDROPS} 等 12 种
\end{itemize}

\paragraph{（3）多标签查询语法。}
\begin{itemize}
    \item 查询方块：\texttt{MATCH (n:item:block) WHERE n.Name CONTAINS '关键词'}
    \item 查询生物：\texttt{MATCH (n:item:monster) WHERE n.Name CONTAINS '关键词'}
    \item 查询纯道具：\texttt{MATCH (n:item) WHERE NOT n:block AND NOT n:monster}
\end{itemize}

\paragraph{（4）路径查询与返回形态。}
本块引导模型优先返回路径，即使用 \texttt{RETURN path} 而非仅返回孤立节点。返回路径时，后端可以直接从记录中提取节点与边，用于绘制关系图；若只返回孤立节点，边信息丢失，前端无法展示实体之间的连接关系。

关于查询方向：无向关系与有向关系在 Cypher 语义上并不等价。在探索性问句中使用无向关系可以扩大命中范围，但也可能匹配到额外的边。本文将此作为提示策略的一个选项，不作为普遍性结论。

\paragraph{（5）少样本示例（few-shot）补充。}
以下标题中“方块挖掘”“生物掉落”等来自本文案例数据，仅作模式演示。

\begin{lstlisting}[language=cypher, caption={方块挖掘（无方向路径）}]
-- 深积岩用什么工具挖
MATCH path = (n:item:block)-[:TOOLMINEDROPS]-(tool:item)
WHERE n.Name CONTAINS '深积岩'
RETURN path LIMIT 10
\end{lstlisting}

\begin{lstlisting}[language=cypher, caption={生物掉落（无方向路径）}]
-- 击败野人会掉落什么
MATCH path = (n:item:monster)-[:DROPS]-(drop:item)
WHERE n.Name CONTAINS '野人'
RETURN path LIMIT 20
\end{lstlisting}

\begin{lstlisting}[language=cypher, caption={配方原料（无方向路径）}]
-- 石剑的合成配方
MATCH path = (product:item)-[:PRODUCES]-(r:recipe)-[:CONSUMES]-(material:item)
WHERE product.Name CONTAINS '石剑'
RETURN path LIMIT 20
\end{lstlisting}

\paragraph{（6）输出格式。}
要求大语言模型返回 JSON，含 \texttt{cql} 与 \texttt{params} 字段。

\subsection{部署与调参中的经验性说明}

在表~\ref{tab:nlq-category} 的自然语言实例测试之外，实现过程中还有以下几点观察。优先使用 \texttt{RETURN path} 与无向关系，有利于表格与力导向关系图同时展示。多标签方案（\texttt{:item:block} 等）使查询路径比旧版桥接关系方案更短，更易于模型理解和生成。\texttt{Name CONTAINS} 适合含复合名称的道具查询。\texttt{\$param} 参数化与关键字黑名单是防止写操作误用的工程措施。以上均为实现过程中的经验性观察，不构成与公开基准数据集的对照评测。

\subsection{迁移场景下的提示词改写说明}

在将系统迁移到其他 Neo4j 知识图谱时，需要改写系统提示词。
现有提示词由两类内容组成，处理方式不同。

第一类是与领域无关的约束性说明，迁移时可以直接保留，无需修改。
具体包括：角色定义（只读 Cypher 生成器，不允许写操作）、输出格式要求（JSON，含 \texttt{cql} 与 \texttt{params} 字段）、路径优先返回策略（使用 \texttt{RETURN path} 代替孤立节点），以及无向关系的使用建议。

第二类是与迷你世界数据库绑定的描述，迁移时需要按目标数据库替换。
具体包括：Schema 描述（第（2）块，含节点标签列表、关系类型列表与各标签的属性键说明）和少样本示例（第（5）块，含示范问句及其对应的 Cypher 语句）。
将这两块的内容替换为目标数据库的对应描述后，其余部分保持不变，即可完成提示词的迁移。

少样本示例的数量与质量对迁移效果影响较大。
建议针对目标数据库各主要节点标签与关系类型各准备一至两个代表性问句，
覆盖单节点属性查询、路径查询和带过滤条件的查询三种形态，使模型对目标领域的典型 Cypher 模式有足够的参照。
 引入
\section{自然语言提示词与调试说明补遗}
\label{app:prompt-supplement}

本附录给出提示词的完整分块内容与少样本示例，供读者对照正文实现或在答辩演示中复现。对应的提示词文件位于仓库 \texttt{backend/app/system\_prompt.txt}，与第~\ref{sec:detailed-implementation}~章第~\ref{subsec:nlq-impl}~节的说明一致；如二者有出入，以仓库当前文件为准。

\subsection{提示词分块}

\paragraph{（1）角色定义。}
告知模型它的任务是根据用户问句生成 Cypher 查询语句，且生成的语句必须是只读的，不允许包含任何写操作。

\paragraph{（2）Schema 描述。}
向模型说明当前图数据库的节点与关系结构：
\begin{itemize}
    \item 基础标签：\texttt{:item}，通用属性含 ID、Name、Type、Disc、GetWay，所有道具节点均带此标签
    \item 方块节点带 \texttt{:item:block} 双标签，在通用属性基础上增加 MineTool、ToolLevel
    \item 生物节点带 \texttt{:item:monster} 双标签，增加 Life、Attack
    \item 其他标签：\texttt{:group}（分组）、\texttt{:recipe}（配方）、\texttt{:smelt}（熔炼）、\texttt{:device}（设备）
    \item 关系类型：\texttt{IN\_GROUP}、\texttt{CONSUMES}、\texttt{PRODUCES}、\texttt{DROPS}、\texttt{TOOLMINEDROPS} 等 12 种
\end{itemize}

\paragraph{（3）多标签查询语法。}
\begin{itemize}
    \item 查询方块：\texttt{MATCH (n:item:block) WHERE n.Name CONTAINS '关键词'}
    \item 查询生物：\texttt{MATCH (n:item:monster) WHERE n.Name CONTAINS '关键词'}
    \item 查询纯道具：\texttt{MATCH (n:item) WHERE NOT n:block AND NOT n:monster}
\end{itemize}

\paragraph{（4）路径查询与返回形态。}
本块引导模型优先返回路径，即使用 \texttt{RETURN path} 而非仅返回孤立节点。返回路径时，后端可以直接从记录中提取节点与边，用于绘制关系图；若只返回孤立节点，边信息丢失，前端无法展示实体之间的连接关系。

关于查询方向：无向关系与有向关系在 Cypher 语义上并不等价。在探索性问句中使用无向关系可以扩大命中范围，但也可能匹配到额外的边。本文将此作为提示策略的一个选项，不作为普遍性结论。

\paragraph{（5）少样本示例（few-shot）补充。}
以下标题中“方块挖掘”“生物掉落”等来自本文案例数据，仅作模式演示。

\begin{lstlisting}[language=cypher, caption={方块挖掘（无方向路径）}]
-- 深积岩用什么工具挖
MATCH path = (n:item:block)-[:TOOLMINEDROPS]-(tool:item)
WHERE n.Name CONTAINS '深积岩'
RETURN path LIMIT 10
\end{lstlisting}

\begin{lstlisting}[language=cypher, caption={生物掉落（无方向路径）}]
-- 击败野人会掉落什么
MATCH path = (n:item:monster)-[:DROPS]-(drop:item)
WHERE n.Name CONTAINS '野人'
RETURN path LIMIT 20
\end{lstlisting}

\begin{lstlisting}[language=cypher, caption={配方原料（无方向路径）}]
-- 石剑的合成配方
MATCH path = (product:item)-[:PRODUCES]-(r:recipe)-[:CONSUMES]-(material:item)
WHERE product.Name CONTAINS '石剑'
RETURN path LIMIT 20
\end{lstlisting}

\paragraph{（6）输出格式。}
要求大语言模型返回 JSON，含 \texttt{cql} 与 \texttt{params} 字段。

\subsection{部署与调参中的经验性说明}

在表~\ref{tab:nlq-category} 的自然语言实例测试之外，实现过程中还有以下几点观察。优先使用 \texttt{RETURN path} 与无向关系，有利于表格与力导向关系图同时展示。多标签方案（\texttt{:item:block} 等）使查询路径比旧版桥接关系方案更短，更易于模型理解和生成。\texttt{Name CONTAINS} 适合含复合名称的道具查询。\texttt{\$param} 参数化与关键字黑名单是防止写操作误用的工程措施。以上均为实现过程中的经验性观察，不构成与公开基准数据集的对照评测。

\subsection{迁移场景下的提示词改写说明}

在将系统迁移到其他 Neo4j 知识图谱时，需要改写系统提示词。
现有提示词由两类内容组成，处理方式不同。

第一类是与领域无关的约束性说明，迁移时可以直接保留，无需修改。
具体包括：角色定义（只读 Cypher 生成器，不允许写操作）、输出格式要求（JSON，含 \texttt{cql} 与 \texttt{params} 字段）、路径优先返回策略（使用 \texttt{RETURN path} 代替孤立节点），以及无向关系的使用建议。

第二类是与迷你世界数据库绑定的描述，迁移时需要按目标数据库替换。
具体包括：Schema 描述（第（2）块，含节点标签列表、关系类型列表与各标签的属性键说明）和少样本示例（第（5）块，含示范问句及其对应的 Cypher 语句）。
将这两块的内容替换为目标数据库的对应描述后，其余部分保持不变，即可完成提示词的迁移。

少样本示例的数量与质量对迁移效果影响较大。
建议针对目标数据库各主要节点标签与关系类型各准备一至两个代表性问句，
覆盖单节点属性查询、路径查询和带过滤条件的查询三种形态，使模型对目标领域的典型 Cypher 模式有足够的参照。
 引入
\section{自然语言提示词与调试说明补遗}
\label{app:prompt-supplement}

本附录给出提示词的完整分块内容与少样本示例，供读者对照正文实现或在答辩演示中复现。对应的提示词文件位于仓库 \texttt{backend/app/system\_prompt.txt}，与第~\ref{sec:detailed-implementation}~章第~\ref{subsec:nlq-impl}~节的说明一致；如二者有出入，以仓库当前文件为准。

\subsection{提示词分块}

\paragraph{（1）角色定义。}
告知模型它的任务是根据用户问句生成 Cypher 查询语句，且生成的语句必须是只读的，不允许包含任何写操作。

\paragraph{（2）Schema 描述。}
向模型说明当前图数据库的节点与关系结构：
\begin{itemize}
    \item 基础标签：\texttt{:item}，通用属性含 ID、Name、Type、Disc、GetWay，所有道具节点均带此标签
    \item 方块节点带 \texttt{:item:block} 双标签，在通用属性基础上增加 MineTool、ToolLevel
    \item 生物节点带 \texttt{:item:monster} 双标签，增加 Life、Attack
    \item 其他标签：\texttt{:group}（分组）、\texttt{:recipe}（配方）、\texttt{:smelt}（熔炼）、\texttt{:device}（设备）
    \item 关系类型：\texttt{IN\_GROUP}、\texttt{CONSUMES}、\texttt{PRODUCES}、\texttt{DROPS}、\texttt{TOOLMINEDROPS} 等 12 种
\end{itemize}

\paragraph{（3）多标签查询语法。}
\begin{itemize}
    \item 查询方块：\texttt{MATCH (n:item:block) WHERE n.Name CONTAINS '关键词'}
    \item 查询生物：\texttt{MATCH (n:item:monster) WHERE n.Name CONTAINS '关键词'}
    \item 查询纯道具：\texttt{MATCH (n:item) WHERE NOT n:block AND NOT n:monster}
\end{itemize}

\paragraph{（4）路径查询与返回形态。}
本块引导模型优先返回路径，即使用 \texttt{RETURN path} 而非仅返回孤立节点。返回路径时，后端可以直接从记录中提取节点与边，用于绘制关系图；若只返回孤立节点，边信息丢失，前端无法展示实体之间的连接关系。

关于查询方向：无向关系与有向关系在 Cypher 语义上并不等价。在探索性问句中使用无向关系可以扩大命中范围，但也可能匹配到额外的边。本文将此作为提示策略的一个选项，不作为普遍性结论。

\paragraph{（5）少样本示例（few-shot）补充。}
以下标题中“方块挖掘”“生物掉落”等来自本文案例数据，仅作模式演示。

\begin{lstlisting}[language=cypher, caption={方块挖掘（无方向路径）}]
-- 深积岩用什么工具挖
MATCH path = (n:item:block)-[:TOOLMINEDROPS]-(tool:item)
WHERE n.Name CONTAINS '深积岩'
RETURN path LIMIT 10
\end{lstlisting}

\begin{lstlisting}[language=cypher, caption={生物掉落（无方向路径）}]
-- 击败野人会掉落什么
MATCH path = (n:item:monster)-[:DROPS]-(drop:item)
WHERE n.Name CONTAINS '野人'
RETURN path LIMIT 20
\end{lstlisting}

\begin{lstlisting}[language=cypher, caption={配方原料（无方向路径）}]
-- 石剑的合成配方
MATCH path = (product:item)-[:PRODUCES]-(r:recipe)-[:CONSUMES]-(material:item)
WHERE product.Name CONTAINS '石剑'
RETURN path LIMIT 20
\end{lstlisting}

\paragraph{（6）输出格式。}
要求大语言模型返回 JSON，含 \texttt{cql} 与 \texttt{params} 字段。

\subsection{部署与调参中的经验性说明}

在表~\ref{tab:nlq-category} 的自然语言实例测试之外，实现过程中还有以下几点观察。优先使用 \texttt{RETURN path} 与无向关系，有利于表格与力导向关系图同时展示。多标签方案（\texttt{:item:block} 等）使查询路径比旧版桥接关系方案更短，更易于模型理解和生成。\texttt{Name CONTAINS} 适合含复合名称的道具查询。\texttt{\$param} 参数化与关键字黑名单是防止写操作误用的工程措施。以上均为实现过程中的经验性观察，不构成与公开基准数据集的对照评测。

\subsection{迁移场景下的提示词改写说明}

在将系统迁移到其他 Neo4j 知识图谱时，需要改写系统提示词。
现有提示词由两类内容组成，处理方式不同。

第一类是与领域无关的约束性说明，迁移时可以直接保留，无需修改。
具体包括：角色定义（只读 Cypher 生成器，不允许写操作）、输出格式要求（JSON，含 \texttt{cql} 与 \texttt{params} 字段）、路径优先返回策略（使用 \texttt{RETURN path} 代替孤立节点），以及无向关系的使用建议。

第二类是与迷你世界数据库绑定的描述，迁移时需要按目标数据库替换。
具体包括：Schema 描述（第（2）块，含节点标签列表、关系类型列表与各标签的属性键说明）和少样本示例（第（5）块，含示范问句及其对应的 Cypher 语句）。
将这两块的内容替换为目标数据库的对应描述后，其余部分保持不变，即可完成提示词的迁移。

少样本示例的数量与质量对迁移效果影响较大。
建议针对目标数据库各主要节点标签与关系类型各准备一至两个代表性问句，
覆盖单节点属性查询、路径查询和带过滤条件的查询三种形态，使模型对目标领域的典型 Cypher 模式有足够的参照。
